\section{Results}
\label{sec:results}

\subsection{Enforcement Costs: Prices, Competition, and Tender Success}
\label{sec:main_results}

Urgency distorts every stage of the procurement process. Reference prices---set during the planning phase---are \BPrefPctPreferred{} higher in the preferred specification ($p < 0.10$; Online Appendix, Table~A.4). The raw premium is \BPrefPctItemFE{} with item FE alone but attenuates sharply with PBU FE, indicating that buyer-level heterogeneity accounts for much of the unconditional gap.

Negotiated prices are the core outcome (Table~\ref{tab:neg_prices}). Panel~A reports the total effect; Panel~B controls for log quantity.

\BPvtab{tab_neg_prices}

\paragraph{Total effect (Panel~A).}
Negotiated prices are significantly higher for urgent purchases across all specifications. The total effect ranges from \BPnegPctItemFE{} ($e^{\BPnegCoefItem}-1$) with item FE only to \BPnegPctRange{} with our preferred specifications (columns~3--4). The attenuation pattern mirrors that of reference prices, confirming the importance of buyer-level heterogeneity.

\paragraph{Mediation decomposition (Panel~B).}
Adding log quantity as a control reduces the coefficient on urgency by roughly half in the most saturated specifications, from \BPnegPanelAcoefShort{} to \BPnegPanelBcoefShort{} (column~3).\label{fn:mediation-anchor} Because quantity is itself an outcome of urgency---urgent orders are by design smaller---Panel~B does not identify a ``direct effect'' in the structural-causal sense: log quantity is a post-treatment variable, and conditioning on it re-weights the comparison toward observations with similar order size \citep{acharya2016explaining}. We therefore read the contrast between Panels~A and~B as a descriptive mediation decomposition: of the \BPnegPctHeadline{} total premium (Panel~A), roughly half co-moves with quantity differences induced by urgency, while a residual \BPnegPanelBpct{} ($e^{\BPnegPanelBcoefShort}-1$, $p < 0.10$) remains once order-size variation is partialled out. The coefficient on log quantity is large and negative ($\BPnegPanelBqty$), consistent with the bulk-discount channel.

\subsection{Competition and Tender Success}
\label{sec:competition}

Urgent purchases attract \BPfirmsPctHeadline{} fewer bidding firms ($\BPfirmsCoef$, $p < 0.01$; Online Appendix, Table~A.6), consistent with reduced competition under compressed timelines \citep{bulow1996auctions}. Quantities also fall, though imprecisely ($\BPqtyCoef$, not significant; Online Appendix, Table~A.5). Yet urgent tenders are \BPsuccessPP{} \emph{more} likely to succeed ($p < 0.01$; Online Appendix, Table~A.7)---precisely what the ``under the gun'' mechanism predicts: officials accept worse terms to ensure the purchase goes through. This also rules out sample selection: urgent tenders succeed \emph{despite} less favorable conditions.

\subsection{The ``Under the Gun'' Effect}
\label{sec:utg}

The sharpest test compares litigated and administrative urgent purchases (Table~\ref{tab:underthegun}). Both share the compressed internal phase, the small-order, diverted-budget regime of urgent procurement; they differ only in whether the officer faces personal fines and fund seizure if the tender fails. The comparison is deliberately \emph{conservative}: administrative requests pass through a scientific committee that screens on evidence-based medicine criteria, so the administrative counterfactual represents the best case the health system achieves under urgency. Any residual premium therefore understates the pure sanction margin for items the committee would have rejected, and overstates it to the extent that litigated cases are systematically more severe or complex within item. We read the \BPutgPctRange{} as a descriptive decomposition of the cost gap between penalty regimes, not a structural estimate of the sanction margin in isolation.

\BPvtab{tab_underthegun}

\paragraph{Total effect (Panel~A).}
Administrative purchases have significantly lower negotiated prices than litigated ones across all specifications. The implied premium ranges from \BPutgPctItem{} with item fixed effects alone (column~1) to \BPutgPctItemYM{} with item and year-month fixed effects plus PBU controls (column~4). In our preferred specification (column~3), the coefficient is $\BPutgCoef$ ($p < 0.05$), implying that litigated purchases are \BPutgPct{} more expensive than comparable administrative purchases.

\paragraph{Mediation decomposition (Panel~B).}
Adding log quantity as a control pulls the coefficient down to \BPutgPanelBcoef{} and strips statistical significance (SE~$= \BPutgPanelBSE$, column~3). Applying the same mediation reading as in \S\ref{sec:main_results}, the \BPutgPctRange{} total cost premium largely co-moves with quantity differences between administrative and litigated purchases---sanctions prevent aggregation, yielding smaller orders that forgo bulk discounts---rather than with a residual per-unit price differential, which is economically modest and statistically indistinguishable from zero. A back-of-envelope calculation is consistent with this reading: the log-quantity gap between administrative and litigated purchases (\BPqtyGapAdmLit) times the bulk-discount elasticity ($\BPbulkElast$) accounts for a price gap of about \BPqtyImpliedPct{} ($e^{\BPqtyGapAdmLit \times |\BPbulkElast|} - 1$) from quantity fragmentation alone, more than twice the \BPutgPct{} total. The fiscal cost of sanctions thus shows up primarily through demand fragmentation, not through officials paying more per comparable unit. Framework agreements for commonly litigated items would address this channel directly.

With item$\times$year-month fixed effects---comparing administrative and litigated purchases of the \emph{same item in the same month}---the coefficient is virtually unchanged ($\BPutgCoefItemYM$, $p < 0.05$, $N = \BPutgNobsItemYM$), ruling out within-item time-varying confounders as a driver.\footnote{Item$\times$year-month FE absorb \BPutgSingletonsItemYM{} singleton observations, leaving \BPutgNobsItemYM{} with within-cell variation. Of \BPutgCellsItemYM{} item$\times$month cells, \BPutgCellsBoth{} (\BPutgCellsBothPct) contain both purchase types. The identifying variation, while limited, produces estimates indistinguishable from the baseline.} Adding log quantity to this saturated specification drives the coefficient close to zero ($\BPutgPanelBItemYMcoef$, SE~$= \BPutgPanelBItemYMSE$), confirming that the premium largely tracks demand fragmentation even under the tightest possible comparison. Regressing log quantity on the administrative indicator under our preferred specification (item, year, and PBU fixed effects) yields a coefficient of $\BPutgQtyAdminCoefPref$: for the same item, administrative orders are roughly \BPutgQtyAdminFoldPref{} times the size of litigated orders. Sanctions do not merely correlate with smaller orders; they prevent demand aggregation by a multiple that, applied to the bulk-discount elasticity above, more than accounts for the entire UTG price gap.

\paragraph{Dose-response.} The mechanism predicts that the urgency premium should scale with the \emph{intensity} of judicial pressure on the item, not with a binary litigation indicator. Binning items by court orders per item-year, the urgency premium rises sharply from \BPdoseLow{} at 1--2 orders (95\%~CI [\BPdoseLowCIlo, \BPdoseLowCIhi]) to \BPdoseMid{} at 3--5 orders ([\BPdoseMidCIlo, \BPdoseMidCIhi]), and remains high at \BPdoseHigh{} for 6--10 orders ([\BPdoseHighCIlo, \BPdoseHighCIhi]).\footnote{The mid and high bins are not statistically distinguishable from each other ($z = \BPdoseMidVsHighZ$), consistent with the saturation pattern; the lo--mid contrast is the binding evidence of a gradient.} The pattern is consistent with a dose-response that saturates: the marginal cost of an additional court order is largest among items at the low end of judicial exposure, and flattens once an item is already heavily litigated. Pure case-mix selection would predict no such gradient.

These results demonstrate that judicial sanctions are a quantitatively important driver of procurement costs. The \BPutgPctRange{} total premium represents a substantial and previously unquantified cost of judicial enforcement in health procurement, operating largely through the quantity channel.

\subsection{Falsification: Items Never Litigated}
\label{sec:placebo}

If the urgency premium reflects the causal effect of judicial pressure
rather than unobserved item characteristics, it should be absent for
items that are never subject to court orders. Running the preferred
specification on items with zero litigated purchases across
\BPsampleStartYear--\BPsampleEndYear{}
(restricted to those with both ordinary and administrative purchases)
yields placebo coefficients of $\BPplaceboNegCoef$ (SE~$= \BPplaceboNegSE$) for negotiated
prices and $\BPplaceboRefCoef$ (SE~$= \BPplaceboRefSE$) for reference prices---both
economically small and statistically insignificant (Online Appendix,
Table~\ref{tab:placebo}). The urgency premium exists only for items
that are targets of litigation, ruling out a generic correlation
between urgency and prices.

\subsection{Demand vs.\ Supply Side}
\label{sec:supplier_fe}

Table~\ref{tab:supplier_fe} decomposes the urgency premium on negotiated
prices by adding supplier fixed effects. Since \BPsupFirmsBothShare{} of winning firms
serve both ordinary and urgent tenders, the within-firm variation
identifies how much of the premium reflects \emph{the same supplier
pricing differently under urgency} rather than \emph{officials selecting
into worse matches} when time-pressed.

\BPotab{tab_supplier_fe}

Column~(1) reproduces the preferred baseline (\BPsupBaselinePct{} premium, $p < 0.01$).
Adding firm fixed effects (column~2) cuts the coefficient by \BPsupFirmFEAttn\% to
$\BPsupFirmFE$ and eliminates statistical significance: roughly half of the total
premium is absorbed by supplier selection under pressure, and the
remaining half is a within-firm markup that the same supplier extracts
when the purchase is court-mandated. With both firm FE and log quantity
(column~3), the residual coefficient falls to $\BPsupFirmFEqty$ (insignificant).
The premium therefore splits into two distinct channels---demand
fragmentation (captured by the quantity margin) and supplier exploitation
(captured by the within-firm markup)---each with a different policy
implication, which we discuss in Section~\ref{sec:conclusion}.

\subsection{Where the Cost Bites Hardest}
\label{sec:heterogeneity}

Splitting the sample by competition intensity, the urgency premium is
three times larger in competitive markets than in concentrated ones
($\BPhetCompetitive$ vs.\ $\BPhetConcentrated$ log points; interaction $\BPhetInteraction$, $p < 0.01$; Online
Appendix, Table~\ref{tab:het_competition}). The direction is initially
counter-intuitive, since standard auction theory would assign more
rent-extraction scope to sellers in concentrated markets
\citep{bulow1996auctions}. Our reading is that in competitive markets
the equilibrium relies on a thin margin of entrants whose participation
is itself compressed by the short internal phase: when planning time
shrinks, the participation margin collapses faster than the price
margin widens, and observed prices rise sharply. In concentrated
markets the premium exists but is partly absorbed by pre-existing
rents. We treat this heterogeneity as descriptive of where the
enforcement cost bites hardest, not as an additional causal test.

\subsection{Summary of Effects}
\label{sec:summary}

Figure~\ref{fig:coefplot} summarizes the paper's quantitative findings
in one panel. Each row plots the estimated coefficient and 95\%
confidence interval from the preferred specification (item + year + PBU
FE, clustering at the PBU level). The pattern that emerges is
economically coherent: urgent purchases raise both reference and
negotiated prices, reduce bidder participation, and nonetheless
\emph{succeed} at higher rates---consistent with officials accepting
worse terms to ensure compliance. The negative coefficient on the
administrative indicator in the ``under the gun'' block confirms that
litigated purchases cost substantially more than administrative ones,
even though both share identical planning constraints.

\begin{figure}[ht]
  \centering
  \includegraphics[width=0.85\textwidth]{fig_08_coefplot.pdf}
  \caption{Coefficient Estimates with 95\% Confidence Intervals
  (Preferred Specification)}
  \label{fig:coefplot}
  \par\smallskip\footnotesize\textit{Notes:} Coefficients from the
  preferred specification (item + year + PBU fixed effects), with
  standard errors clustered at the PBU level. The first block shows the
  effect of urgency (administrative + litigated, relative to ordinary)
  on five outcomes: reference prices, negotiated prices, log number of
  bidders, and tender success (linear probability). The second block
  reports the ``under the gun'' coefficient---log negotiated price
  regressed on an administrative indicator within the urgent subsample.
\end{figure}
